\chapter[Natural Selection Is Not Evolution by Effort]{Natural Selection\protect\\ Is Not Evolution by Effort}

\begin{kaichang}
Open the medicine cabinet at home, and you may well find a pack of antibiotics: amoxicillin, a cephalosporin, or something else. The leaflet will almost certainly contain a sentence like ``Complete the full course even if your symptoms improve.'' Doctors repeat the warning: even when your fever and sore throat are gone, do not stop the medicine yourself.

That seems strange. If you are better, why keep taking medicine? Some people guess that drug companies want to sell more pills; others think the infection might not be completely gone and could return. Both guesses sound somewhat reasonable, but neither is the real answer. That answer lies in a process of ``selection'' that has operated for billions of years and may be unfolding inside a patient right now. Write down your guess, and we will return to it at the chapter's end.

A larger question also awaits. Many people explain evolution this way: giraffes have long necks because their ancestors ``kept stretching to reach high leaves, generation after generation, until their necks grew longer.'' In other words, organisms can evolve in the direction they need simply by trying hard. Almost everyone has entertained this thought---and almost every part of it is wrong. Where does it go wrong? Let us start with those antibiotics.
\end{kaichang}

\textbf{Translator safety note:} Take antibiotics only as prescribed, and ask the prescriber or pharmacist before changing treatment. Do not use leftover or someone else's antibiotics. This chapter's general explanation of treatment duration is oversimplified; clinicians select an evidence-based effective course for the particular infection. See \href{https://www.cdc.gov/antibiotic-use/about/}{CDC antibiotic guidance} and \href{https://www.nice.org.uk/guidance/ng15/evidence/full-guideline-pdf-252320799}{NICE antimicrobial stewardship guidance}.

\section{An Invention Losing Its Effectiveness}

First consider some observable, measurable facts.

In 1928, Fleming noticed that bacteria surrounding a Penicillium mold in a dish had died, leading him to discover penicillin. By the 1940s, mass production made previously fatal battlefield wound infections treatable, and deaths from pneumonia and puerperal fever fell sharply. Antibiotics rank among the twentieth century's greatest medical inventions and made a tangible contribution to rising global life expectancy.

Decades later, however, things changed. The 1960s saw the appearance of methicillin-resistant Staphylococcus aureus, or MRSA, a ``superbug'' able to survive on hospital sheets and door handles and withstand ordinary antibiotics. Tuberculosis treatment also changed: ordinary TB can mostly be cured with six months of medicine, whereas multidrug-resistant TB requires eighteen months or longer, with worse side effects and lower cure rates. The World Health Organization estimates that over a million deaths worldwide in 2019 were directly related to bacterial resistance. This is a model estimate, not a person-by-person count, but the scale is alarming enough.

\textbf{Translator safety note:} The source's tuberculosis-duration comparison is not current treatment guidance: WHO now recommends shorter regimens, including six-month options, for eligible patients with drug-resistant TB. Diagnosis, drug selection, and duration require specialist assessment; do not alter treatment from this example. See \href{https://tbksp.who.int/en/node/3032}{WHO drug-resistant TB treatment guidance}.

Here is the key question: not a single atom in the antibiotic molecule has changed during those decades. Penicillin remains the same molecule and kills by the same mechanism. \textbf{The drug has not deteriorated; the bacteria have changed.} Bacteria do not hold planning meetings or resolve to work harder. How do they ``change''?

\section{Bacteria Already Differ}

Zoom down to the particle level. As Chapter 72 explained, a bacterium's genetic material is a circular DNA molecule that must be copied before division. Copying is precise but imperfect: base pairing occasionally goes wrong, and although repair catches most mistakes, some slip through. Those are mutations.

Get a feel for the scale. Under suitable conditions, an E. coli bacterium can reproduce about every 20 minutes. Roughly $10^{13}$ bacteria live in your gut. Each cell acquires, on average, about $10^{-3}$ new mutations per replication. That individual number seems small, but multiply it by astronomical population sizes and rapid reproduction, and the conclusion reverses: \textbf{in any enormous bacterial population, almost every possible single-point mutation is occurring in some bacterium right now}.

Most mutations are neutral: a changed base does not affect protein function, or lies at a noncritical position, and the bacterium carries on as usual. Some are harmful, slowing growth or preventing survival. A tiny minority happen to block an antibiotic: some change the shape of its protein target so the drug molecule no longer fits; others add a ``pump'' to the cell wall that expels incoming drug molecules. Only one or two bacteria among ten million companions may have such changes.

Now the antibiotic arrives. It kills $99.9\%$ of the bacteria, including every individual lacking that mutation. The one or two bacteria that happen to carry resistance mutations survive, with newly vacant territory and nutrients before them. They copy themselves and divide: twenty minutes later, two become four, then eight. Within days, a new population arises in which \textbf{every member descends from a survivor and is resistant from birth}.

Look carefully at this machine's logic: bacteria do not make mutations ``in order to'' resist drugs. Mutations arise independently of the drug---they occurred for billions of years before antibiotics appeared, and soil bacteria carried resistance genes long before humans used penicillin. The antibiotic does just one thing: \textbf{kills the ill-suited and leaves those that happen to fit}. It supplies no direction; it is simply a sieve.

\begin{zhengju}
Do mutations arise before the drug appears, or does the drug ``induce'' them? In 1943, Luria and Delbrück answered with an elegant experiment. They divided the same E. coli stock among many small culture tubes, then spread each tube's bacteria on separate dishes covered with bacteriophages, viruses that infect bacteria. They counted the surviving phage-resistant colonies from each tube.

If phages induced resistance, every bacterium would have the same probability of becoming resistant on encountering a phage. Colony counts should therefore be similar across tubes, like the uniform distribution of die rolls. If resistance mutations instead occurred randomly, early or late, before phage exposure, the outcome would differ completely. A ``lucky'' tube with an early resistance mutation would contain generations of descendants, thousands or tens of thousands by plating time, and might produce hundreds of colonies; most tubes without an early mutation would produce few. The latter pattern appeared: colony counts varied enormously, producing a ``jackpot'' distribution. Mutation preceded selection, and randomness preceded need. This combination of thought experiment and experimental design later earned them a Nobel Prize and drove the first nail into the evolutionary principle ``mutation is undirected; selection is directional.''
\end{zhengju}

\textbf{Translator safety note:} The bacterial-culture experiment is historical evidence, not a home activity. Do not culture unknown microbes or select antibiotic-resistant organisms; use published data or the paper simulation below. See \href{https://institute.acs.org/acs-center/lab-safety/education-training/safer-experiments.html}{ACS risk-assessment guidance}.

\section{Three Conditions, One Machine}

We can now dismantle the machine. Natural selection needs just three conditions to hold together:

\begin{itemize}[itemsep=2pt]
\item \textbf{Heritable differences}: individuals in a population differ, and DNA can pass those differences to offspring. Chapters 72 and 73 explained where differences arise.
\item \textbf{Differences in survival and reproduction}: in the \emph{current} environment, certain differences help their bearers live longer and leave more offspring. The key is ``reproductive success,'' not merely staying alive. A long life with no offspring counts as zero in evolution's ledger.
\item \textbf{Enough time}: generation after generation, bearers of advantageous differences become a larger fraction of the population.
\end{itemize}

All three are purely mechanical conditions. None mentions ``desire,'' ``effort,'' or ``progress.'' Differences arise randomly, the environment automatically selects, and the rest is arithmetic: whoever makes more copies has a louder voice in the next generation.

To clarify ``a louder voice,'' we need a term. Different versions of a gene are \textbf{alleles}, as Chapter 71 explained: resistance is one version, susceptibility another. An allele's proportion among all copies in a population is its \textbf{allele frequency}. Evolution can be defined precisely in one sentence: \textbf{changes in a population's allele frequencies across generations}. When frequencies change, evolution has occurred, with no necessary connection to ``progress'' or being ``more advanced.''

\begin{qiaoliang}
How fast can this machine run? Suppose an infection contains one million bacteria, exactly one of which carries a resistance mutation: a frequency of $10^{-6}$, or one in a million. An antibiotic kills $99.9\%$ of susceptible bacteria, leaving about 1,000 of the original 999,999. The resistant bacterium survives.

The survivors then reproduce in the vacant territory. Suppose that after one treatment course, before the immune system finishes clearing the infection, each surviving bacterium leaves an average of 1,000 descendants. Susceptible survivors become about $10^6$, and resistant ones about $10^3$. Wait---resistant bacteria still seem to be a minority! Remember round two: another antibiotic round, perhaps after transmission to another person, cuts the $10^6$ susceptible bacteria to $10^3$, while the $10^3$ resistant bacteria remain unharmed and expand to $10^6$. After two rounds, the resistance allele frequency has risen from one in a million to about $99.9\%$.

From almost zero to almost complete dominance in just two selection rounds: this is the arithmetic doctors fear most, and direct evidence that evolution can be fast. It need not take millions of years; if differences are consequential enough and reproduction fast enough, a few days suffice.
\end{qiaoliang}

\section{Moths, Beaks, and Changing Winds}

Bacteria are hidden in culture dishes, out of everyday sight and touch. Fortunately, natural selection's three conditions apply to all organisms, including visible animals. Two classic cases deserve a closer look.

\subsection{Peppered Moths: Frequencies Written on Bark}

Britain's peppered moth was almost exclusively pale in early nineteenth-century collection records. Its whitish-gray, black-speckled wings made it difficult for birds to spot against pale, lichen-covered bark. Around 1848, a dark individual was first recorded near Manchester; by 1895, dark moths exceeded $90\%$ locally. In fifty years, an allele had gone from rarity to dominance. Why Manchester, a city dense with Industrial Revolution chimneys?

The causal chain runs as follows: coal pollution killed pale lichens and blackened trunks. Pale moths on dark bark went from ``invisible'' to ``conspicuous,'' increasing their risk of bird predation; dark moths gained the opposite advantage. Notice the causal direction: soot did not ``blacken the moths.'' An individual moth's color stayed unchanged throughout its life. What changed was \emph{each color's opportunity to leave offspring within the population}.

\begin{zhengju}
Was this explanation correct? Could soot directly change moth genes instead? In the 1950s, ecologist Kettlewell tested the idea by releasing marked moths of both colors in polluted countryside and unpolluted woodland, then attempting to recapture them. Dark moths had significantly higher recapture rates in polluted areas, with the reverse in clean areas. He also directly observed birds taking conspicuous individuals.

Critics later identified weaknesses: naturally, peppered moths mostly rest high in the canopy rather than on the trunks where he placed them, potentially exaggerating the difference. This is science working as it should: peers challenge experimental design, and conclusions temporarily acquire question marks. In the 2000s, Majerus repeated the work more naturally, releasing moths at their actual resting positions and observing for years. Again, bird predation discriminated by color strongly enough to explain the frequency changes. Meanwhile, after Britain controlled air pollution, bark became lighter and dark-moth frequency indeed fell. \textbf{The prediction came true}: reverse the environment, and selection reverses direction. A theory truly captures causality when it can predict the direction of future change.
\end{zhengju}

\subsection{Darwin's Finches: Reversal Before Our Eyes}

If the moth story unfolded over fifty years, Galápagos ground finches compress evolution into a few. In 1977, the Grants were measuring medium ground-finch beaks annually on the islands. A severe drought exhausted small, easily opened seeds first, leaving large, hard ones. Many birds with small, weak beaks could not crack them and starved; birds with slightly deeper, stronger beaks survived and reproduced. Chicks hatched the next year had an average beak depth about $4\%$ greater than before the drought---only a few parts in a hundred, but a clear, measurable shift.

The story did not end with ``ever-larger beaks.'' In 1982--1983, El Niño brought torrential rain, grasses flourished, and small seeds became abundant again. Now large, clumsy beaks were disadvantaged because they handled small seeds inefficiently. After several seasons, average beak depth fell again. \textbf{The same population was pushed one way by selection and then back the other way within a few years.} Nothing is eternally ``better''; there is only ``more advantageous here and now.'' Darwin saw static differences between islands' beaks; his successors saw living, reversible evolution on a single island.

\subsection{Back to the Medicine Bottle}

Antibiotic resistance is the same machine's third example and the closest to your own life. Like finch beaks, it contains clues to reversal: many resistance mutations have costs. The drug-expelling ``pump'' consumes energy, and a reshaped protein may work more slowly. Without antibiotics, resistant bacteria often lose to efficient susceptible bacteria, and their frequency gradually declines. That is why antibiotic misuse---taking them for colds, stopping after two days, or feeding them to livestock as growth promoters---harms everyone: it keeps the ``sieve'' constantly suspended over bacteria, retaining resistant types generation after generation.

\section{What Changes, and What Does Not}

We can now dismantle the most common misconception. Compare the three cases, and the same pattern appears:

\begin{itemize}[itemsep=2pt]
\item No individual moth ``became'' dark; dark moths' \emph{frequency} in the population changed.
\item No finch ``lengthened'' its beak during the drought; deeper-beaked birds left more offspring, shifting the next generation's average.
\item No bacterium ``trained itself'' to resist drugs; individuals that happened to resist survived and monopolized the next generation.
\end{itemize}

An individual's lifetime follows \textbf{development}: a fertilized egg grows, matures, and ages according to gene-regulatory networks (Chapter 75), with its genotype largely unchanged. Evolution occurs not within individuals but between generations of \textbf{populations}, through changing allele frequencies. \textbf{Individuals develop; populations evolve.} Confusing these is the root of the ``effort produces evolution'' illusion.

The giraffe story also needs retelling. Ancestors did not ``stretch their necks hard and pass stretched necks to their children.'' Instead, heritable differences in neck length already existed within the herd. In seasons with scarce low vegetation, slightly longer-necked individuals reached more leaves and left more offspring. Over many generations, the population's average neck length shifted. No giraffe had to try, and none needed to: selection automatically did the work.

\begin{wuqu}
``Use and disuse; inheritance of acquired characteristics.'' This was Lamarck's famous proposal before Darwin: a muscular blacksmith would have sons born with muscular arms, and giraffes stretching their necks would produce longer-necked offspring. Its mistake was treating \emph{individual changes} as \emph{heritable changes}. However large your trained muscles become, germ-cell DNA does not rewrite itself to match them. Weightlifters' children are born like other children and must train for themselves. Counterexamples are easy to find: families that have worn earrings for generations do not produce baby girls with pierced earlobes.

In fairness, Chapter 74 showed that some epigenetic states can pass across generations, and environment can indeed affect developing traits, as Chapter 73 explained. But such phenomena are not ``directional''---organisms cannot customize changes on demand---and do not overturn the main account: DNA-level changes remain the fundamental source of heritable variation, and frequency changes still depend on population-level selection. Lamarck guessed the wrong mechanism; Darwin and later genetics supplied the right one.
\end{wuqu}

\section{Adaptation Answers Only to Here and Now}

The three cases teach another lesson: \textbf{adaptation does not point toward the future; it settles accounts with the current environment}. Dark moths excel during pollution and fall behind when the air clears; a large beak saves lives in drought but becomes a burden in wet years; resistant bacteria rule a drug-filled world and yield territory to susceptible bacteria when the drug disappears. Evolution has no permanent winners because the environment defines ``winning,'' and environments keep changing.

Thus ``adapted'' does not mean ``perfect,'' only ``good enough now and slightly better than competitors.'' Natural selection is a make-do engineer: it does not design from scratch but chooses among existing variants and accepts what works. Vertebrate eyes have a blind spot where the optic nerve passes through the retina without photoreceptors; the crossing of human air and food passages in the throat makes choking possible. These ``design flaws'' are not records of evolution failing but its signature: only a process of patching old plans until they work well enough would leave such traces. An engineer designing from scratch would not work this way.

Evolution therefore has neither endpoint nor ladder. It does not climb from ``lower'' to ``higher,'' much less toward humanity. Chapter 78 will explore the tree's actual shape. Bacteria have undergone selection for billions of years and are still bacteria, superbly adapted to their own worlds. Humans are not evolution's purpose, merely one of countless branches.

\begin{shiyan}
\textbf{Simulation: Be a ``Bird'' and Do the Selecting}

Materials: cover a table with intricately patterned wrapping paper or a patterned tablecloth. Prepare 50 paper pieces of each of two colors, one close to the background and one strongly contrasting, plus tweezers and a timer.

Procedure: (1) Mix and scatter both colors on the cloth. (2) Have a classmate act as ``predator,'' picking up pieces individually with tweezers for 20 seconds and setting each aside. (3) Count how many of each color were taken. (4) Let ``survivors'' reproduce: add one matching piece per remaining piece, restoring the total to 100. (5) Repeat steps (2)--(4) for three or four rounds, recording color proportions each time.

What to observe: which color's proportion increases round by round? If your ``predation'' is random enough and the environment, the cloth, stays unchanged, frequency will move consistently in one direction: you have just carried out natural selection. Repeat with a different background. Does the direction change?

Safety reminder: cut the pieces with children's scissors, do not point tweezers at anyone, and keep paper pieces out of mouths.
\end{shiyan}

\section{Back to the Antibiotic Pack}

Now to fulfill the opening promise: why must you finish the course after symptoms disappear?

Because bacterial susceptibility is not uniform but a continuous spectrum: the most susceptible die first, while slightly more resistant individuals last longer. After the first few days, your fever and sore throat subside because the main bacterial force has been destroyed and numbers are low enough for the immune system to handle easily. Yet the survivors are precisely the population's \emph{most resistant minority}. Stopping now removes the sieve halfway through selection. These ``preliminary-round winners'' gain breathing room to reproduce, establishing a population with higher average resistance. You may relapse and also transmit these upgraded bacteria to someone else.

A complete course maintains drug exposure long enough to finish selection and, together with immunity, reduce survivors to zero. \textbf{``Stopping early'' is not saving medicine; it is personally training more resistant bacteria.} Of course, whether an illness needs antibiotics, which drug to use, and for how long are decisions for a doctor. Antibiotics are useless against viral colds and merely select among your resident gut bacteria. Those are matters of medical advice; this book only explains the evolutionary arithmetic behind them.

\textbf{Translator safety note:} The preceding explanation is not a universal account of antibiotic resistance or a reason to prolong treatment. Treatment duration and resistance have a complex relationship; stewardship guidance asks clinicians to use the shortest effective course for the infection. Follow your current prescription and contact your clinician or pharmacist about symptoms, side effects, or duration rather than stopping or extending treatment yourself. See \href{https://www.nice.org.uk/guidance/ng15/evidence/full-guideline-pdf-252320799}{NICE guidance} and \href{https://www.cdc.gov/antibiotic-use/about/}{CDC patient advice}.

The same arithmetic operates in fields. A new insecticide often works strikingly well at first but becomes less effective after several years. The chemical has not ``gone bad''; selection has produced a resistant pest population. DDT's history illustrates a complete cycle: in the 1940s it seemed a miracle mosquito killer, yet within a decade resistance frequencies in multiple mosquito species went from rare to predominant. Modern agricultural countermeasures explicitly use evolutionary thinking: rotate chemicals with different mechanisms to keep changing the ``sieve's'' direction, or leave an unsprayed ``refuge'' beside a field so susceptible individuals survive and continually dilute resistance alleles. Once you understand that ``selection merely filters,'' these strategies become deductions rather than memorized rules.

\section{When the Machine Falters, and What Comes Next}

Natural selection is powerful, but remember the conditions under which it cannot operate:

\begin{itemize}[itemsep=2pt]
\item Without relevant heritable differences, even strong selection has nothing to work with. Diversity itself is therefore a population's survival capital.
\item If differences affect neither survival nor reproduction, selection is ``blind'' to them. Their allele frequencies can still change through sheer luck: genetic drift.
\item In small populations, luck can even overwhelm advantage: an excellent allele may disappear because its bearer never got to reproduce.
\end{itemize}

The last two cases are not patches added to this chapter's account; they reveal another machine. Evolution is more than natural selection. Mutation supplies material, drift rolls dice, and gene flow carries alleles between populations. Chapter 77 explores these ``nonselective'' forces and how they turn molecular change into a readable ``clock.'' When mountains or oceans separate a population and each part accumulates changes, how do two species eventually form? That is Chapter 78's story.

\begin{tiaozhan}
(You may skip this without losing the main thread.) Can we express selection-driven frequency change mathematically? Yes. Let an allele have frequency $p$, with its bearers leaving an average of $w_1$ offspring per generation, called \textbf{fitness}; bearers of the other version have fitness $w_2$. Set relative fitnesses $w_1 = 1$ and $w_2 = 1 - s$, where $s$ is the \textbf{selection coefficient}: $s = 0.1$ means that version leaves $10\%$ fewer offspring per generation on average. After one generation, frequency approximately becomes
\[
p' = \frac{p}{p + (1-p)(1-s)}.
\]
For resistant bacteria, substitute $p = 10^{-6}$ and a susceptible-type $s$ near $0.999$ in the drug's presence. What is $p'$ after one generation? Then consider peppered moths: estimates put the dark type's advantage at an $s$ on the order of $0.3$. A seemingly modest advantage can raise frequency from $1\%$ to over $90\%$ in fifty generations, or fifty years for moths with one generation per year. Verify this with a calculator: evolution is neither fast nor slow; it runs exactly like compound interest.
\end{tiaozhan}

\section*{Try It Yourself}

\begin{sikaoti}
\item Using ``individuals develop; populations evolve,'' rewrite this sentence and identify its error: ``This finch tried hard to grow a larger beak during the drought to adapt to hard seeds.''
\item Draw an information-flow diagram from ``DNA copying error'' to ``increasing population resistance frequency.'' Identify the random and directional steps and explain what supplies the direction.
\item Without antibiotics, resistant bacteria often gradually decline in frequency. Explain using adaptation to the current environment and resistance costs. Predict what happens if antibiotics return.
\item A report says, ``Bacteria evolved resistance in order to fight antibiotics.'' Identify at least three misleading expressions and rewrite the statement accurately.
\item Suppose a new predator introduced to an island eats only insects in flight. Which traits might change in frequency after several generations? What three conditions does your prediction require?
\item Refer to the compound-interest arithmetic in the quantitative bridge box: an allele's frequency does not double every generation, but how does ``advantageous types leave more offspring'' resemble compound interest on bank savings? Where does the analogy break down?
\end{sikaoti}
